\documentclass[11pt]{article}
\usepackage[margin=1in]{geometry}
\usepackage{amsmath}
\usepackage{booktabs}
\usepackage{url}
\title{Target parameters from the estimator's direct output}
\date{Prepared October 5, 2026}
\begin{document}
\maketitle

\section{Question}
In the joint estimator--controller design (commit \texttt{06af0a2},
\path{JOINT_DESIGN_20261005.tex}), a window fit replaced the tracker's target
scalar acceleration $A$ and sideslip $\beta$ (item E2). The fit used the last
2~s of radar detections. The user asked for the estimator's direct output
instead. This report records the change and what it does to the noisy-sensor
encounters.

\section{Change}
The runtime again publishes the tracker's own reconstruction of $A$ and
$\beta$. The tracker state is $[\rho;q;s]$: relative position, target velocity
and target acceleration in the ego frame. From it,
$A=q^\top s/|q|$ and $\Omega=(Jq)^\top s/|q|^2$
(\path{estimator/nrmmTargetTrackerDerivative.m}). The following were removed:
\begin{itemize}
\item \path{estimator/nrmmTargetParameterFit.m};
\item its calls in \path{estimator/onlineNrmmTrackingRuntime.m}: reset,
  initialization, per-detection history, the override of the published $A$ and
  $\beta$, and the widening of the published radii by the tracker--fit distance;
\item the configuration field \texttt{parameterFitMinimumDuration};
\item the corresponding shift terms in \path{estimator/nrmmControllerErrorBounds.m}.
  This file and \path{config/nrmmTrackingConfig.m} are again identical to
  \texttt{300b70b}.
\end{itemize}
Four unit tests of the fit were removed with it. Everything else in
\texttt{06af0a2} is unchanged: the force-balance ego measurement and the
computed domain (E1), and acquisition and initialization in the adapter (E3).
The controller changes C1--C6 are unchanged too.

\section{Results}
\paragraph{Encounters.} The 98 runs used the same configuration as in
\path{JOINT_DESIGN_20261005.tex}: 14 encounters with exact states, and the
same 14 with noisy sensors for seeds 20261003 and 1--5. Each ran up to 2000
holds (100~s).
\begin{center}
\begin{tabular}{lrrrrrr}
\toprule
 & & & no & & min.\ & min.\ road\\
 & runs & recovered & solution & collision & gap (m) & margin (m)\\
\midrule
exact states & 14 & 14 & 0 & 0 & 0.153 & 3.85\\
noisy, window fit & 84 & 84 & 0 & 0 & 0.416 & 1.00\\
noisy, tracker output & 84 & 82 & 2 & 0 & 0.362 & 0.23\\
\bottomrule
\end{tabular}
\end{center}
The exact-state runs do not use the estimator. Their results are identical to
\texttt{06af0a2}. Seeds 20261003, 1, 2 and 4 recovered 14 of 14; seeds 3 and 5
recovered 13 of 14. No run collided or left the road.

The other noisy runs changed systematically. In the 15-m/s curved runs the
smallest road margin fell to 0.23--0.89~m, from 1.89--3.30~m with the fit. The
minimum gap to the target grew to 3.5--3.9~m, from 1.8--3.2~m. Faced with the
noisier forecast, the controller passed the target more widely and came closer
to the road edge. The median minimum gap over all noisy runs was 1.66~m
(fit: 1.64~m). On the straight road, recovery was confirmed after a median of
10.9~s (fit: 10.9~s). The maximum was 16.7~s, in an 8-m/s braking lead (fit:
13.5~s). The ego estimation errors were unchanged within run-to-run variation:
95th-percentile heading 0.0019--0.0068~rad, lateral velocity
0.010--0.114~m/s, position 0.007--0.022~m
(\path{ESTIMATOR_DIRECT_TARGET_20261005/campaign-metrics.csv}).

\paragraph{Target parameters.} The table gives the root-mean-square error of
the published parameter in each of the 84 noisy runs, as the median (maximum)
over runs
(\path{ESTIMATOR_DIRECT_TARGET_20261005/parameter-errors.csv}).
\begin{center}
\begin{tabular}{lrr}
\toprule
 & window fit & tracker output\\
\midrule
$A$, 0.5--3~s (m/s$^2$) & 0.019 (0.048) & 0.269 (0.461)\\
$\beta$, 0.5--3~s (rad) & 0.0005 (0.0013) & 0.0078 (0.0269)\\
$A$, first 0.5~s (m/s$^2$) & 0.228 (0.550) & 0.543 (1.163)\\
$\beta$, first 0.5~s (rad) & 0.0000 (0.0077) & 0.0208 (0.0498)\\
\bottomrule
\end{tabular}
\end{center}
The tracker's $A$ also changes strongly from one 50-ms frame to the next. It
can sit at the contract limit $\pm1.1$~m/s$^2$ for several frames (below).

\paragraph{The two runs without a solution.} Both stopped in a frame whose plan
needed a long horizon, while the tracker's $A$ was far from the truth
(\path{ESTIMATOR_DIRECT_TARGET_20261005/failures.csv}).
\begin{itemize}
\item \emph{Seed 3, 8-m/s braking lead, $t=2.65$~s.} The lead brakes at
  $A=-0.5$~m/s$^2$. In the eight preceding frames the published $A$ was
  $-0.717$, $-0.898$, $-1.100$, $-1.100$, $-1.100$, $-1.100$, $-0.585$ and
  $-0.594$~m/s$^2$; in the failing frame it was $-0.078$, with
  $\beta=0.008$ instead of 0. The horizon was 84 holds (4.2~s). At its end the
  forecast lead was 3.7~m further ahead and 1.6~m further left than the
  true one.
\item \emph{Seed 5, 8-m/s curved crossing, $t=7.10$~s.} The target turns with
  $A=0.5$~m/s$^2$ and $\beta=0.040$. In the eight preceding frames the
  published $A$ ranged from 0.16 to 0.73~m/s$^2$; in the failing frame it was
  0.78. Until then the horizon had been 8 holds, and the separating speed of
  the plan's endpoint had fallen from 1.83 to 0.56~m/s, close to the
  rollout's 0.5-m/s margin. In the failing frame the rollout ran on to 104
  holds (5.2~s). At the end of that horizon the forecast target was 5.6~m from
  the true one.
\end{itemize}
The failing frames were re-solved with groups of rows removed. Without the
collision rows, or without the terminal separation rows, both problems were
solved. Removing the road, handling or road-recovery rows did not help. With
the collision rows made soft, the slack was nonzero only at the end of the
horizon: nodes 58--85 of 84 holds, and 102--105 of 104. The conflict lies
between the forecast target near the end of a long horizon and the terminal
separation requirement. Beyond the first hold the forecast uses the published
$A$ and $\beta$ without an enclosure (see the scope in
\path{JOINT_DESIGN_20261005.tex}), so the error of $A$ enters it
quadratically in time.

\paragraph{Horizon length.} The horizon is variable and is unchanged by this
change. The shifted plan is rolled out until its first node in the terminal
set, at most 512 holds (25.6~s). Over all noisy frames the median was 16 holds
(0.8~s). The 90th percentile was 22, the 99th 79 and the maximum 138 holds.
87.8\% of frames used at most 16 holds, 2.0\% more than 64 and 0.017\% more
than 128. With the fit the figures were 16, 23, 89 and 135
(\path{ESTIMATOR_DIRECT_TARGET_20261005/horizon-steps.csv}). The longest
horizons occur in the 8-m/s braking lead (138), the 8-m/s curved crossing
(129) and the 15-m/s braking lead (127). The two failures fall in the first
two of these.

\paragraph{Truth audit.} \path{scripts/replayEstimatorValidity.m} replayed
the 14 seed-20261003 runs and reproduced the published estimates to
$5.3\times10^{-15}$. The ego bound was valid in all 5,958 frames, and neither
the ego bound nor the target bound (2,291 frames) failed to contain the truth.
The true state left the shared domain in three frames, through the
rear-adhesion ratio: 1.0072 and 1.0053 at 1.25~s in the 15-m/s head-on and
accelerating head-on, and 1.0004 at 3.75~s in the 15-m/s braking lead. The
sideslip cone held. The target prediction-parameter set is as wide as before.
Per-run medians were a course radius of 0.24--0.53~rad and a speed half-width
of 1.9--4.0~m/s, with the whole contract range for $A$ and curvature
(\path{ESTIMATOR_DIRECT_TARGET_20261005/audit.csv}).

\paragraph{Tests.} 653 of 653 tests pass: the 657 of \texttt{06af0a2} less
the four fit tests (\path{ESTIMATOR_DIRECT_TARGET_20261005/tests.csv}).

\section{Scope}
\begin{itemize}
\item The certified bounds are unaffected: the fit had only replaced the point
  values of $A$ and $\beta$ and widened the radii by its distance to the
  tracker. The failures are not bound violations. They come from the forecast
  beyond the first hold, which no certificate covers.
\item The tracker's gains minimize a bound on its output-noise velocity
  variance subject to the Lipschitz dissipation inequality. The noise of the
  acceleration state $s$, from which $A$ and $\beta$ are reconstructed, is not
  part of that design. Ways to keep the direct output and reduce this noise
  include a bound on the variance of $s$ in the gain program, or a
  constant-parameter stage inside the estimator cascade. On the controller
  side, the long horizons could trust the forecast less. None of these was
  tried here.
\item The campaign ran as eight parallel processes; the computation time was
  not re-measured.
\end{itemize}

\section{Addendum: decision}
The remedies of Section~4 were studied in
\path{TRACKER_NOISE_AND_DESIGN_ABLATIONS_20261005.tex}. On October 5, 2026
the user accepted 82 of 84 as the current result.

\section{Reproduction}
The runner is the one of \path{JOINT_DESIGN_20261005.tex}; a copy is in
\path{ESTIMATOR_DIRECT_TARGET_20261005/campaignOne.m.txt}. For one seed:
\begin{verbatim}
c = estimatorControllerIntegrationConfig(); c.randomSeed = seed;
runPotentialFieldCampaign(root, out, 2000, EstimatorConfiguration=c)
\end{verbatim}
The truth audit is \texttt{replayEstimatorValidity(out)} on the
seed-20261003 outputs, and the tests are \texttt{runtests('tests')}. The
source hashes of the files that produced these results are listed in
\path{ESTIMATOR_DIRECT_TARGET_20261005/source-sha256.txt}. The outputs are
under\\
\path{~/.cache/collisionAvoidance/estimator-direct-20261005}.
\end{document}
